\chapter{Tests, déploiement et résultats}
\label{ch:tests}

\section{Stratégie de test}

Les scénarios suivent \url{https://lexhub.space} : inscription OTP puis login JWT ; atelier en cinq étapes ; assistant limité aux plaidoiries \textbf{validées} ou \textbf{finalisées} ; isolation entre cabinets ; packs masqués pendant l'essai ; invitation e-mail ; popups de quotas ; export Word (.docx) et PDF ; factures listées uniquement si montant $> 0$~MAD.

\begin{table}[H]
\centering
\caption{Scénarios de validation (extrait)}
\label{tab:scenarios}
\small
\renewcommand{\arraystretch}{1.2}
\begin{tabular}{@{}p{3.6cm}p{5.2cm}p{3.6cm}@{}}
\toprule
\textbf{Scénario} & \textbf{Attendu} & \textbf{Résultat} \\
\midrule
Inscription + OTP & Cabinet actif, connexion & Conforme \\
Brouillon & Invisible assistant & Conforme \\
Valider pour l'assistant & Visible de tous les assistants & Conforme \\
Assigner & Validé + mise en avant & Conforme \\
Clarification QCM & 0--5 questions puis rédaction & Conforme \\
Assistant révision (chat) & Réponse + document mis à jour & Conforme \\
Micro requête (Terminé) & Whisper auto FR/AR & Conforme \\
Export Word (.docx) / PDF & RTL + armoiries + en-tête ; finalisée & Conforme \\
Isolation Rabat / Atlas & Aucune donnée croisée & Conforme \\
Essai : pas de packs & Offre masquée & Conforme \\
Repère IA du jour & Popup continuer / demain & Conforme \\
Forfait essai / mois épuisé & Popup de guidance & Conforme \\
Multi-support mobile ($\le 880$\,px) & Dock fixé, rôles OK (avocat / assistant / admin) & Conforme \\
\bottomrule
\end{tabular}
\end{table}

\section{Déploiement et intégration continue (CI/CD)}
\label{sec:cicd}

La production s'appuie sur Caddy, Docker Compose et Alembic. GitHub Actions publie les images \texttt{frontend}, \texttt{backend} et \texttt{celery} vers GHCR. Cette chaîne est livrée.

\section{Résultats obtenus}

L'atelier en cinq étapes, la clarification QCM, l'assistant de révision (réponse + édition), l'export officiel, la règle « validé = visible par l'assistant », l'isolation multi-tenant, l'\textbf{expérience multi-support} (desktop et mobile, section~\ref{sec:mobile}) et le modèle économique du chapitre~\ref{ch:economie} sont en production. Le paiement marchand réel n'est pas encore activé.

\section{Risques et limites}
\label{sec:limites}

\begin{enumerate}[leftmargin=1.4em]
  \item Un micro faible dégrade la transcription : la relecture (marques orange) reste obligatoire.
  \item La dictée et la rédaction dépendent d'APIs cloud.
  \item La correction de forme n'est pas une vérification juridique.
  \item LexHub n'intègre pas Mahakim et n'est pas un ERP de facturation.
  \item L'avocat reste responsable du contenu.
\end{enumerate}

\section{Perspectives}
\label{sec:perspectives}

\begin{enumerate}[leftmargin=1.4em]
  \item Activer l'\textbf{encaissement réel} (carte opérateur) sur le socle de paiement déjà préparé.
  \item Améliorer la \textbf{qualité STT} et les aides de relecture pour les conditions acoustiques difficiles.
  \item Renforcer la continuité \textbf{dossier $\rightarrow$ audience} (rappels, clé de barre, exports).
  \item Préparer des \textbf{interopérabilités} (calendriers, imports de pièces) sans casser les maîtres de donnée MADI.
\end{enumerate}
